\documentclass[11pt,a4paper]{article}

\usepackage[left=2cm,text={17cm,24cm},top=3cm]{geometry}
\usepackage[english]{babel}
\usepackage[utf8]{inputenc}
\usepackage[T1]{fontenc}

\usepackage{url}
\usepackage{float}
\usepackage{amsthm}
\usepackage{comment}
\usepackage{etoolbox}
\usepackage{graphicx}
\usepackage{hyperref}
\usepackage{tocloft}

\def\UrlBreaks{\do\/\do-\do\&\do=\do\_\do?} % URL breaking characters

\newcommand{\red}[1]{\textcolor{red}{#1}} % \red{text in red}
\newcommand{\blue}[1]{\textcolor{blue}{#1}} % \blue{text in blue}
\newcommand{\TODO}{\textbf{\textcolor{red}{TODO}}} % red bold TODO
\newcommand{\tilda}{\raisebox{0.5ex}{\texttildelow}} % command \tilda for '~' character

\renewcommand{\cftdot}{} % remove dots ([sub]secion ......... number) from table of contents

\theoremstyle{definition}
\newtheorem{definition}{Definition}[section]
\graphicspath{{img/}} % path to images
\patchcmd{\thebibliography}{\section*{\refname}}{}{}{} % do not create section for bibliography
\hypersetup{
    linktoc    = all,
    colorlinks = true,
    citecolor  = green,
    linkcolor  = red,
    urlcolor   = blue,
}

\begin{document}

\begin{titlepage}

    \begin{center}
        % FIX: lines must end with '%', if not then it will result in an incorrect centering
        \vfill {%
            \Huge{%
                \textsc{%
                    Faculty of Informatics\\[3mm]%
                    Masaryk University%
                }%
            }%
        }%

        \hfill\\[15mm]

        \begin{figure}[!h]
            \centering
            \includegraphics[scale=3]{img/muni-fi-logo.pdf}
        \end{figure}

        \hfill\\[10mm]

        \Huge{
            \textbf{
                IV064
            }
        }

        \hfill\\[-10mm]

        \huge{
            \textbf{
                Information Society
            }
        }

        \hfill\\[10mm]

        \LARGE{
            \textbf{
                Open access to source code, and results, open systems
            }
        }
        \vfill

    \end{center}

        \Large{
            \noindent Adrián Tóth (491322)\hfill \today
        }

\end{titlepage}

\setlength{\parskip}{0pt}
    {
        \hypersetup{
            hidelinks=true
        }
        \tableofcontents
    }
\setlength{\parskip}{0pt}

\newpage

\section{The Open Access}

    Nowadays, the open access term is used in many ways. Firstly, let us declare the meaning of this expression within the context of IT to better understand other things that will be mentioned later. You may heard phrases like '\textit{Open access to source code}' as well as '\textit{Open source software}'. Accordingly, this phrases may also appear under the '\textit{free}' acronym, which is not the proper notion. There is an article~\cite{WP:free-vs-open-source} about the contrast of these two terms, which describes and explains the differences between them. All of these names indicates one attribute (model, philosophy or methodology) of the software development - the right of free access to the source code of a software to anyone known as the open source philosophy. This right stands for the free way of software source code free usage, inspection, modification and distribution that may be restricted by further open source licenses such as GPL~--~GNU General Public License, which will be described later. To clarify the veritable meaning of the open source let us specify the correct definition itself - the open source definition.\\[-3mm]

    \begin{definition}
        \cite{BOOK:open-source-def, WP:opensource-osd, WP:opensource-debian}\\[-5mm]
        \begin{center}
            \begin{minipage}{0.9\textwidth}
                Open source does not just mean access to the source code. The distribution terms of an open source software must comply with the following criteria:\\[-7mm]
                \begin{itemize}
                    \item Free Redistribution\\[-7mm]
                    \item Inclusion of Source Code\\[-7mm]
                    \item Inclusion of Derived Works\\[-7mm]
                    \item Integrity of The Author's Source Code\\[-7mm]
                    \item No Discrimination Against Persons or Groups\\[-7mm]
                    \item No Discrimination Against Fields of Endeavor\\[-7mm]
                    \item Distribution of License\\[-7mm]
                    \item License Must Not Be Specific to a Product\\[-7mm]
                    \item License Must Not Restrict Other Software\\[-7mm]
                    \item License Must Be Technology--Neutral\\
                \end{itemize}
            \end{minipage}
        \end{center}
    \end{definition}

    As has been noted above, the definition is denoting a specific software type which is made freely available with respect to modification and distribution to anyone including license restrictions and non-discriminatory rules. At the present time, this influencing fact has affected not just the software itself but the development as well. Open source projects, products, or initiatives embrace and celebrate principles of open exchange, collaborative participation, rapid prototyping, transparency, meritocracy, and community-oriented development~\cite{WP:what-is-os}. The strongest side of the open source is its community together with the collaboration power. The open source community is diverse and highly motivated~\cite{WP:opensource-com}.\\

    At the present time, there is a countless amount of open source projects. Many of them become popular quite quickly and has drawn the interest of academia and industry. Several projects are running for decades, some of them are just at the beginning of their dawn. As an illustration, the Table~\ref{fig:open-source-projects} below contains some of the well known and still running projects under the terms of open source.\\

    \begin{table}[H]
        \begin{center}
            \begin{tabular}{l|l|l}
                Project      & Year      & Author(s)                              \\
                \hline
                Unix OS      & 1969/1970 & Ken Thompson, Dennis Ritchie \& others \\
                GNU Project  & 1983      & Richard Stallman                       \\
                Linux Kernel & 1991      & Linus Torvalds                         \\[-5mm]
            \end{tabular}
        \end{center}
        \caption{The oldest and still existing open source projects~\cite{SD:linux-kernel, WP:gnu-proj-init, DLACM:history-of-unix}.}
        \label{fig:open-source-projects}
    \end{table}

    In contrast to the open source, its antonym is known as proprietary sometimes referred as non-free. In conclusion, it indicates copyright restrictions that prevent unrestricted distribution or reuse of the software. The pros and cons of these two types of software philosophy will be explained later in the following sections.

\section{Success of Open Source Projects}

    In the face of the fact, that nowadays it is very hard to be successful without any reward for the given effort, open source projects have their characteristic attribute that they are invincible in this field. The roots of success lie in the kindness and helpfulness of people who are developing software for the others as a hobby and for free. These developers, representing and formulating the community, are the crucial heroes for the success.

    \subsection{The Open Source Community}

    \subsection{Collaboration and Contribution}

    \subsection{Source Code Availability and Hosting}

    \subsection{Feedback and Support}

\section{Open Source Consequences}

    \subsection{Influencing Impact of Open Source}

    \subsection{Threats, Weaknesses and Vulnerabilities}

    \subsection{Pros and Cons}

\section{The Open Systems}

    \TODO

\section{My Notes}

    \TODO REMOVE!

    \begin{itemize}
      \item \url{https://ieeexplore.ieee.org/document/4556977}\\[-8mm]
      \item \url{https://ieeexplore.ieee.org/document/5662568}\\[-8mm]
      \item \url{https://ieeexplore.ieee.org/document/1620054}\\[-8mm]
      \item \url{https://www.sciencedirect.com/science/article/pii/B9781555583200500027}\\[-8mm]
      \item \url{https://www.sciencedirect.com/science/article/pii/B9781555583200500167}\\[-8mm]
    \end{itemize}

\newpage

\section{References}
\bibliographystyle{englishiso}
\begin{flushleft}
    \bibliography{quotation}
\end{flushleft}

\end{document}
